% Options for packages loaded elsewhere
\PassOptionsToPackage{unicode}{hyperref}
\PassOptionsToPackage{hyphens}{url}
\documentclass[
  11pt,
  a4paper,
]{article}
\usepackage{xcolor}
\usepackage[margin=1in]{geometry}
\usepackage{amsmath,amssymb}
\setcounter{secnumdepth}{-\maxdimen} % remove section numbering
\usepackage{iftex}
\ifPDFTeX
  \usepackage[T1]{fontenc}
  \usepackage[utf8]{inputenc}
  \usepackage{textcomp} % provide euro and other symbols
\else % if luatex or xetex
  \usepackage{unicode-math} % this also loads fontspec
  \defaultfontfeatures{Scale=MatchLowercase}
  \defaultfontfeatures[\rmfamily]{Ligatures=TeX,Scale=1}
\fi
\usepackage{lmodern}
\ifPDFTeX\else
  % xetex/luatex font selection
  \setmainfont[]{STIXGeneral}
  \setmonofont[]{Menlo}
\fi
% Use upquote if available, for straight quotes in verbatim environments
\IfFileExists{upquote.sty}{\usepackage{upquote}}{}
\IfFileExists{microtype.sty}{% use microtype if available
  \usepackage[]{microtype}
  \UseMicrotypeSet[protrusion]{basicmath} % disable protrusion for tt fonts
}{}
\makeatletter
\@ifundefined{KOMAClassName}{% if non-KOMA class
  \IfFileExists{parskip.sty}{%
    \usepackage{parskip}
  }{% else
    \setlength{\parindent}{0pt}
    \setlength{\parskip}{6pt plus 2pt minus 1pt}}
}{% if KOMA class
  \KOMAoptions{parskip=half}}
\makeatother
\usepackage{longtable,booktabs,array}
\newcounter{none} % for unnumbered tables
\usepackage{calc} % for calculating minipage widths
% Correct order of tables after \paragraph or \subparagraph
\usepackage{etoolbox}
\makeatletter
\patchcmd\longtable{\par}{\if@noskipsec\mbox{}\fi\par}{}{}
\makeatother
% Allow footnotes in longtable head/foot
\IfFileExists{footnotehyper.sty}{\usepackage{footnotehyper}}{\usepackage{footnote}}
\makesavenoteenv{longtable}
\setlength{\emergencystretch}{3em} % prevent overfull lines
\providecommand{\tightlist}{%
  \setlength{\itemsep}{0pt}\setlength{\parskip}{0pt}}
\usepackage{newunicodechar}
\newfontfamily\numeralfallback{STIX Two Text}
\newunicodechar{∎}{\ensuremath{\blacksquare}}
\usepackage{fvextra}
\DefineVerbatimEnvironment{Highlighting}{Verbatim}{breaklines,breakanywhere,breakanywheresymbolpre={},breakanywheresymbolpost={},breaksymbolleft={},breaksymbolright={},commandchars=\\\{\}}
\newunicodechar{⁰}{{\numeralfallback\char"2070\relax}}
\newunicodechar{₀}{{\numeralfallback\char"2080\relax}}
\newunicodechar{₁}{{\numeralfallback\char"2081\relax}}
\newunicodechar{₂}{{\numeralfallback\char"2082\relax}}
\newunicodechar{₃}{{\numeralfallback\char"2083\relax}}
\newunicodechar{₄}{{\numeralfallback\char"2084\relax}}

\usepackage{seqsplit}
\renewcommand{\texttt}[1]{{\ttfamily\seqsplit{#1}}}
\setlength{\emergencystretch}{6em}

\newunicodechar{⁻}{{\numeralfallback\char"207B\relax}}
\newunicodechar{⁸}{{\numeralfallback\char"2078\relax}}
\newunicodechar{ⁱ}{{\numeralfallback\char"2071\relax}}
\newunicodechar{⁴}{{\numeralfallback\char"2074\relax}}
\newunicodechar{⁷}{{\numeralfallback\char"2077\relax}}
\newunicodechar{⁶}{{\numeralfallback\char"2076\relax}}
\newunicodechar{⁵}{{\numeralfallback\char"2075\relax}}
\newunicodechar{₈}{{\numeralfallback\char"2088\relax}}
\newunicodechar{ᵢ}{\textsubscript{i}}
\newunicodechar{ⱼ}{\textsubscript{j}}
\newunicodechar{ₜ}{\textsubscript{t}}
\usepackage{bookmark}
\IfFileExists{xurl.sty}{\usepackage{xurl}}{} % add URL line breaks if available
\urlstyle{same}
% fallback for those not using the hyperref driver hyperxmp:
\makeatletter
\@ifundefined{xmpquote}{\newcommand{\xmpquote}[1]{#1}}{}
\makeatother
\hypersetup{
  hidelinks,
  pdfcreator={LaTeX via pandoc}}

\author{}
\date{}

\begin{document}

\section{Module C: Halving and Petrial operations on regular
geometries}\label{module-c-halving-and-petrial-operations-on-regular-geometries}

\subsection{Abstract}\label{abstract}

We study two operations on regular incidence geometries through their
distinguished chamber generators and the objects retained by the
operation. Replacing an endpoint involution by the conjugate of its
neighbour has an exact intersection criterion: one overlap equality and
an intersection-closed family of supports. Its geometric realization is
a component of a bipartite double cover. This yields the halving theorem
for regular lattice polytopes, identifies the parity groups of
generalized cubes, and gives their spherical, hyperbolic and toroidal
applications. The odd-endpoint case produces six rank-four geometries
for the first Janko group; an exhaustive proper-residue census proves
completeness and excludes all higher ranks. For geometric Petrials, the
retained vertices and edges instead determine a possibly nonfaithful
action. Its exact kernel leads to the thin-Petrial obstruction and to
faithful regularity and biduality, without a finiteness assumption. The
shared action and coset recognition results are used from the companion
paper \emph{Marked groups, kernels and intersections}.

\subsection{1. Introduction}\label{introduction}

A local change of chamber generators can alter an incidence geometry in
two distinct ways. It may select an index-two subgroup while changing
the parabolic subgroups that represent its objects. It may also identify
objects that become indistinguishable after part of the incidence data
is forgotten. The first issue is controlled by subgroup intersections
and flag completion. The second requires the kernel of the induced
action to be kept explicitly.

The endpoint operation replaces an involution a by aba, leaving its
neighbour b and the other markers unchanged. Section 3 gives a necessary
and sufficient criterion for the new marking to have the intersection
property. An odd endpoint order forces the generated group to remain
unchanged; a proper halving has index two. Both are instances of the
same criterion. Section 4 realizes this change by doubling a graph and
choosing a component, with an exact parity condition on its higher
shadows. This proves the halving theorem for lattice polytopes.

For generalized cubes, the halving group is the even-weight part of a
binary permutation module. Its simple coordinate description gives all
group orders, polygon parameters and locally spherical cases (§5). An
explicit Gray map then identifies the actual toroidal face posets,
affine markings and translation lattices (§6). These results recover the
principal construction and applications of Montero and Ivić Weiss
{[}MW{]}. The same endpoint theorem supplies the finite constructions of
Hartley and Leemans {[}HL{]}. Section 7 recovers all six markings, their
geometries and the exhaustive rank classification; its two complete
finite certificates use a subgroup catalogue, not an earlier list of
geometries.

The Petrial operation of Bonheure, Buekenhout and Leemans {[}BBL{]} has
a different local rule. Faces are actual vertex--edge sets traced by a
chamber walk, and equal sets are identified. Section 8 proves that the
action kernel is the core of the original face-changing involution
subgroup. Thinness then forces a commuting pair of generators, makes the
induced action faithful and regular, and allows the original geometry to
be recovered by a second Petrial. The argument applies to infinite
groups and residues. The source's unconditional embedding assertion must
therefore be replaced by its exact kernel statement; an explicit example
shows the distinction.

Theorems U.Action, U.Geometry, U.Helly and Proposition U.ExtensionOrbits
of the companion paper {[}U{]} supply the common action and coset
interfaces. All hypotheses specific to these operations are established
here. Danzer's generalized-cube construction, classical Coxeter
structure and the complete Janko subgroup catalogue remain named inputs.
The original construction results retain the attribution above; the
finite certificate interface and precise source qualifications are
stated at the end.

\subsection{2. Chamber actions and finite coset
tests}\label{chamber-actions-and-finite-coset-tests}

A geometry has a finite set of types; its flags are sets of pairwise
incident objects of distinct types, every flag extends to a chamber, and
its residues are the induced geometries on objects incident with a fixed
flag. It is thin when every panel has exactly two completions and
residually connected when every residue of rank at least two has
connected incidence graph. Regularity means chamber transitivity of the
full type-preserving automorphism group. No finiteness of the object set
is implicit.

A marking of a group G is an indexed generating family of nonidentity
involutions \((r_i)_{i\in I}\). Write \(G_J=\langle r_j:j\in J\rangle\).
The intersection property is \[
G_J\cap G_K=G_{J\cap K}\qquad(J,K\subseteq I).
\tag{C1}
\] The maximal parabolics are \(P_i=G_{I\setminus\{i\}}\); their right
cosets are incident when they intersect. Typed coset Helly means that
every pairwise intersecting family with distinct types has a common
element. U.Geometry shows that IP and typed Helly give exactly the
regular thin residually connected coset geometry, with G as its full
type-preserving group. IP alone is not sufficient for that conclusion.

\textbf{Lemma C.Gallery (the marked stabilizers).} In a thin regular
residually connected geometry, the full group acts freely on chambers.
At a base chamber let \(r_i\) be the unique group element carrying it
across its i-panel. Then \(r_i\) is an involution, the \(r_i\) generate
the full group, and the stabilizer of a base subflag F is \[
G_F=\langle r_i:i\notin t(F)\rangle.
\tag{C2}
\]

\emph{Proof.} U.Geometry proves chamber-graph connectedness and freeness
of the chamber action. The square of \(r_i\) fixes the base chamber, so
\(r_i\) has order two. A chamber gallery expresses every chamber image
by successive panel changes, which gives generation. Apply the same
gallery argument in the residue at F. Only the missing types can change,
so its chamber-transitive stabilizer is generated by their \(r_i\).
Conversely those \(r_i\) fix F. This proves (C2). ∎

\textbf{Lemma C.HellyCount (finite recognition).} For a finite group,
put \[
X_{ij}=\{P_i x:x\in P_j\},\qquad P_J=\bigcap_{j\in J}P_j.
\] Typed coset Helly is equivalent to \[
\left|\bigcap_{j\in J}X_{ij}\right|=[P_J:P_J\cap P_i]
\quad(i\notin J,\ |J|\ge2).
\tag{C3}
\]

\emph{Proof.} The cosets represented in \(P_J\) form a subset of the
displayed intersection, of size the right side. Equality is therefore
equality of these sets. This is exactly the anchored condition of
U.Helly: a coset meeting every base \(P_j\) has a representative in
their common subgroup. The finite count proves both directions. ∎

For an operation retaining some old objects and encoding the new ones by
their actual subsets, U.Action identifies the kernel with the pointwise
kernel on the retained objects. Sections 4 and 8 will specify the sets,
encodings and equivariant surjections explicitly. Objects that merely
receive the same name or label do not satisfy that theorem's
injective-encoding hypothesis.

\subsection{3. Endpoint conjugation and central
twists}\label{endpoint-conjugation-and-central-twists}

Let \((G;a,b,(r_i)_{i\in E})\) be a C-group marking, with \([a,r_i]=1\)
for all \(i\in E\). No string diagram, finite group, or underlying
geometry is assumed. Define

\[
a'=aba,\quad R_J=\langle r_j:j\in J\rangle,\quad
U_J=\langle a,b,R_J\rangle,\quad H_J=\langle a',b,R_J\rangle,
\] \[
P=\langle b,R_E\rangle,\quad H=H_E,\quad
\mathcal F=\{J\subseteq E:a\in H_J\}.
\tag{E1}
\]

Conjugation by a exchanges b and \(a'\) and fixes \(R_J\) pointwise.
Consequently

\[
H_J\trianglelefteq U_J=H_J\langle a\rangle,
\qquad [U_J:H_J]\le2.
\tag{E2}
\]

\textbf{Theorem C.Endpoint.} The new marking \((H;a',b,(r_i))\) has IP
if and only if

\[
P\cap P^a=R_E,
\qquad \mathcal F\text{ is closed under intersections}.
\tag{E3}
\]

The group H is normal of index at most two in G. It is G precisely when
\(\mathcal F\) is nonempty. Since E is finite, a nonempty
intersection-closed \(\mathcal F\) is precisely the family of supersets
of one least support \(S\subseteq E\).

\textbf{Proof.} The family \(\mathcal F\) is upward closed. Any new
marked subgroup is of one of the forms

\[
A_{\epsilon,J}=\langle b^\epsilon,R_J\rangle\ (\epsilon=0,1),
\qquad H_J,\qquad Q_J^a\text{ with }Q_J=\langle b,R_J\rangle,
\tag{E4}
\]

where \(b^{0}\) means that b is omitted. These correspond to omitting
\(a'\), using both endpoint labels, or using \(a'\) but omitting b.
Intersections of two A's are old marked intersections. Also

\[
A_{\epsilon,I}\cap H_J
\le A_{\epsilon,I}\cap U_J=A_{\epsilon,I\cap J};
\]

all generators on the right belong to \(H_J\), so equality holds.
Intersections of two \(Q^a\)'s follow by conjugating old IP. Conjugation
by a normalizes \(H_I\), so an H-with-\(Q^a\) intersection reduces to
the just proved H-with-A case. For an A-with-\(Q^a\) intersection, the
first condition in (E3) confines it to \(R_E\). Intersecting there and
using that a centralizes \(R_E\) gives \(R_{I\cap J}\).

It remains to intersect \(H_I\) and \(H_J\). Put \(K=I\cap J\). Old IP
gives

\[
H_K\le H_I\cap H_J\le U_I\cap U_J=U_K.
\tag{E5}
\]

If \(a\in H_K\) then \(H_K=U_K\). Otherwise \(H_K\) has index two in
\(U_K\), so the middle subgroup is either \(H_K\) or \(U_K\). The latter
case occurs exactly when a belongs to both \(H_I\) and \(H_J\), and
intersection closure excludes it. This proves all new marked
intersections. Necessity of the first condition follows by intersecting
the two endpoint-deleted new parabolics; necessity of the second follows
by intersecting \(H_I,H_J\). The generators are involutions and the
proved IP supplies independence. Equation (E2) with \(J=E\) gives the
group and index assertions. A finite nonempty upward intersection-closed
family contains its total intersection S, giving the least-support
description. ∎

\textbf{Corollary C.Odd (odd endpoint and proper halving).} If
\(m=|ab|\) is odd, the only test in (E3) is the overlap
\(P\cap P^a=R_E\), and \(H=G\). If H is proper in G, the only test is
again the overlap, and \([G:H]=2\). In either case the endpoint product
has order m/gcd(m,2) for finite m; infinite m remains infinite.

Indeed \(a'b=(ab)^{2}\), and for odd m,

\[
a=(a'b)^{(m+1)/2}b\in H_\varnothing.
\]

Thus \(\mathcal F\) is respectively all subsets or empty. An even
endpoint order by itself does not imply that H is proper; the
least-support alternative in Theorem C.Endpoint cannot be discarded.

\textbf{Lemma C.Twist (central twist).} Suppose b,c are involutions with
\(|bc|=m\) odd and \(m\ge 3\), and w is a nonidentity involution
centralizing b,c.~Replacing b by wb preserves the generated group and
changes \textbar bc\textbar{} to 2m, provided w belonged to the original
group.

\textbf{Proof.} The commuting factors w and bc have coprime orders, so
\((wbc)^m=w\) and \(|wbc|=2m\). Hence the new generators recover w and
b. The new marker wb is a nonidentity involution: \(wb=1\) would make
\(b=w\) commute with c, contrary to \(m\ge 3\) odd. ∎

A parabolic containing w,b,c is unchanged by this twist; one omitting b
is unchanged. Other intersections and typed Helly still require proof.
Lemma C.Twist is the auxiliary operation needed in the finite
application below, not a blanket geometric preservation theorem.

\subsection{4. Graph halving and lattice
polytopes}\label{graph-halving-and-lattice-polytopes}

Let \(\Gamma\) be a finite-rank residually connected incidence geometry,
of rank \(n\ge 3\). Designate types 0 and 1 as vertices and edges.
Assume their truncation is a simple graph Q, and that an edge is
incident with a higher object exactly when both its endpoints are
incident with that object. For a flag F of higher objects, let \(Q_F\)
be its common vertex-edge shadow. All \(Q_F\) are connected and contain
an edge.

U.Geometry supplies the connected chamber graph of every residually
connected finite-rank geometry. Restricting a chamber gallery to any set
of at least two types shows that the corresponding truncation is
connected, because successive complete incidence subgraphs overlap.
Apply this inside the residue at F; chamber extension supplies an edge.
The edge-shadow hypothesis says \(Q_F\) is the induced graph on its
common vertices.

For a connected graph Q, its bipartite double cover D(Q) has vertices
\((v,\varepsilon )\), \(\varepsilon \in \{0,1\}\), with
\((v,\varepsilon )\) adjacent to \((w,1-\varepsilon )\) when v is
adjacent to w. It is connected when Q has an odd cycle, and has two
components otherwise: parity of walks proves both assertions. Choose a
component C of D(Q). Define \(L_C(\Gamma )\) as follows. Its first two
types are (v,0) and (v,1) in C. A higher object is a pair \((x,C_x)\),
where \(C_x\) is a component of \(D(Q_x)\) contained in C. Incidence of
the first two types is adjacency, incidence with a higher object is
membership in \(C_x\), and two higher objects are incident when their
old objects are incident and their chosen components meet.

The precise compatibility requirement is

\[
\text{if every }Q_x\ (x\in F)\text{ is nonbipartite, then }Q_F
\text{ is nonbipartite},
\tag{PC}
\]

for every nonempty higher flag F. It is automatic if the higher shadows
of every flag are nested, or if Q is bipartite.

\textbf{Theorem C.Graph.} The system \(L_C(\Gamma )\) is always a
geometry. It is residually connected exactly when (PC) holds. It is thin
exactly when \(\Gamma\) is thin. A chamber-transitive automorphism group
of \(\Gamma\) acts faithfully and chamber-transitively on
\(L_C(\Gamma )\) when Q is nonbipartite; when Q is bipartite its
index-two color-preserving subgroup does so.

\textbf{Proof.} For a flag \(\Phi =\{(x,C_x):x\in F\}\) of higher lifts,
write \(W=\cap C_x\). If any \(Q_x\) is bipartite, restriction of
\(C_x\) to \(D(Q_F)\) is one component. Any two such restrictions agree:
incidence makes them meet over the connected \(Q_{x,y}\), and hence
select the same parity class there and over \(Q_F\). Nonbipartite
constituents impose no restriction. Thus W is one component if any
constituent is bipartite, and all of \(D(Q_F)\) otherwise. In particular
W always contains an edge, and is connected exactly under (PC) for this
flag. For F empty use \(W=C\).

An edge of W and a chamber of \(\Gamma\) extending its base edge and F
give a chamber extending \(\Phi\). Every additional higher object has a
unique component containing the chosen lift of that edge. The same
extension works for a flag already containing one lifted vertex. If it
contains both, their base edge is unique by simplicity and incident with
all the old higher objects by the shadow condition. This proves the
geometry axiom.

The residue at (v,0) is the old residue \(\Gamma _v\): a neighbour (w,1)
corresponds to the unique edge \{v,w\}, and the component for an
incident higher object is uniquely specified by (v,0). Simplicity and
the shadow condition preserve all incidences. The other color is
identical. Hence residues of flags containing a lifted vertex are
connected. For a flag of higher objects only, its two-type truncation is
W; when W is connected, every other residual object attaches to it
through a chamber. This proves residual connectivity under (PC). If (PC)
fails, choose its higher flag; all constituents have unique lifts but W
has two components. Were the resulting geometry residually connected,
the truncation lemma inside its residue would make W connected. This
proves necessity without treating a truncation as a residue.

Panels containing both new vertex types correspond to old panels missing
the same higher type. A panel missing one new vertex type corresponds to
an old panel missing the edge type, at the retained endpoint. Old panels
missing the vertex type already have two objects by simplicity. These
correspondences prove both directions of thinness.

If Q is nonbipartite, an old chamber
\((v,\{v,w\},x_{2},\ldots ,x_{n-1})\) corresponds bijectively to the new
chamber with (v,0),(w,1) and the components containing them. Old
automorphisms act naturally on this bijection. If Q is bipartite, the
same bijection uses the old chambers whose chosen vertex is in one fixed
color class. A chamber-transitive group contains a color-exchanging
element, since it exchanges the endpoints of a base edge through a
chamber move. Its color-preserving subgroup therefore has index two and
is transitive on the selected chambers, with the same chamber bijection.
Faithfulness in both cases follows from the retained-object argument
below. ∎

\subsubsection{4.1. The retained-object
action}\label{the-retained-object-action}

For completeness, apply U.Action to the precise object sets in this
construction. Let K be G when Q is nonbipartite and its color-preserving
subgroup otherwise. Take X to be the disjoint union of old vertices and
higher objects, and Y the object set of the halved geometry. The map \[
(v,\varepsilon)\mapsto v,\qquad (x,C_x)\mapsto x
\tag{C4}
\] is an equivariant surjection \(Y\to X\). For nonbipartite Q the
chosen C is all D(Q); for bipartite Q it contains exactly one lift of
every old vertex. In the latter case its intersection with the double
cover of each connected shadow \(Q_x\) is one of that double cover's two
components. Thus every higher object also has a lift.

Encode a new vertex by its fixed color tag and old vertex. Encode
\((x,C_x)\) by the old object x and the two subsets
\(\{v:(v,0)\in C_x\}\), \(\{v:(v,1)\in C_x\}\) of X, retaining its old
type as a tag. Components are determined by their vertex sets, so the
encoding is injective and equivariant. U.Action gives \[
\ker(K\curvearrowright Y)=\ker(K\curvearrowright X)=1.
\tag{C5}
\] The last equality follows because an old edge is uniquely determined
by its two endpoints in the simple graph Q. Fixing X therefore fixes all
old objects. When \(\Gamma\) is thin and (PC) holds, the proved geometry
is thin and residually connected, and its K-action is
chamber-transitive. U.Geometry then identifies K with the full
type-preserving group. This argument does not infer the geometry or its
flag properties merely from faithfulness.

Suppose now \(\Gamma\) is regular and thin, with vertex-changing marker
a, edge-changing marker b, and higher markers \(r_i\); assume a commutes
with the higher markers. Under (PC), Theorem C.Graph produces a regular
thin residually connected geometry. Directly on a base edge, its panel
involutions are \(b^a,b,r_i\): b changes one endpoint, aba changes the
other, and each \(r_i\) changes its higher object. U.Geometry identifies
their group as the full type-preserving group H. Lemma C.Gallery
identifies the stabilizer of every base subflag with the subgroup
generated by its missing-type markers, giving IP and identifying the
actual coset incidence with \(L_C(\Gamma )\).

There is also an exact correspondence between the two parity conditions.
In the base residue with group \(U_J\) from (E1), the vertex stabilizer
is \(Q_J=\langle b,R_J\rangle\). Its vertex graph is bipartite precisely
when \(a\notin H_J\). If \(H_J\) is proper, its two cosets give a
coloring, since \(Q_J\le H_J\) and a crosses colors. Conversely the
action on the two colors supplies a character killing b and all \(r_j\)
and sending a to the nontrivial element, so it excludes a from \(H_J\).
In a globally nonbipartite geometry, higher flags reduce by transitivity
to base flags with shadows \(U_{E\setminus S}\). Condition (PC) says

\[
[E\setminus\{i\}\in\mathcal F\text{ for all }i\in S]
\Longrightarrow E\setminus S\in\mathcal F.
\tag{G1}
\]

An upward family containing E satisfies (G1) exactly when it is
principal: let T be the set of i whose complementary coatom belongs to
\(\mathcal F\); (G1) gives \(E\setminus T\in \mathcal F\), and every
member of \(\mathcal F\) contains E\textbackslash T. Thus geometric
parity compatibility is the least-support condition of Theorem
C.Endpoint. The mixed overlap in that theorem has a different role: it
says the two adjacent vertices determine one edge. Indeed P and \(P^a\)
stabilize the endpoints; their intersection is the edge's pointwise
stabilizer \(R_E\) exactly when no other edge has those endpoints, using
transitivity on panels and the regular chamber identification.

\subsubsection{4.2. Lattice-polytope
halvings}\label{lattice-polytope-halvings}

\textbf{Corollary C.Halving.} Let P be a regular rank-n abstract
polytope, \(n\ge 3\), whose face poset is a lattice. Replacing one
endpoint marker a by aba gives a C-group H normal of index at most two
in Aut(P). Its maximal-parabolic coset system is a regular hypertope,
with full type-preserving group H. The diagram duplicates the edge
labels from b to the remaining markers, and the \(a'–b\) label is
m/gcd(m,2) for \(m=|ab|\) finite, or infinity for infinite m. The index
is two exactly when the corresponding vertex graph is bipartite; for a
final-end operation use the dual graph on facets.

\textbf{Proof.} In a lattice polytope an edge is the join of its two
distinct vertices. Thus two vertices determine at most one edge, and
that edge belongs to a higher face exactly when both vertices do. Higher
faces of a flag form a chain, so their shadows are nested and (PC) is
automatic. Theorem C.Graph and the panel-marker identification prove the
asserted geometry, full group, and IP without an
alternating-semiregular-polytope theorem. Equation (E2) and the coloring
argument give the index. Finally \(a'b=(ab)^{2}\) and a commutes with
all \(r_i\), so \(a'r_i\) is conjugate to \(br_i\); all other products
are unchanged. ∎

This is {[}MW, Theorem 3.1, Proposition 3.2 and Corollary 3.3{]}, after
reversing their final-end convention. Their term \emph{non-degenerate}
is used in the lattice-face-poset sense, as explicitly explained by the
later primary halving manuscript cited in the references; a general
abstract polytope with simple diagram is not substituted for that
hypothesis. N013's auxiliary direct-product Lemma 2.4 follows at once
because a marked subgroup of a product with disjoint factor markings is
the product of the corresponding marked subgroups, and intersections are
coordinatewise. Its quoted Propositions 2.1 and 2.3 are background; the
latter is not needed. U.Geometry supplies the content needed from N013's
quoted Theorem 2.2.

\subsection{5. Parity modules and locally spherical
geometries}\label{parity-modules-and-locally-spherical-geometries}

Retain the following Danzer construction. If K is a finite regular
lattice polytope with m vertices, \(2^K\) has binary vertex set
\(V=F_{2}^m\). A face associated with a face F of K and \(x\in V\)
allows exactly the coordinates indexed by the vertices of F to vary.
These coordinate subcubes, ordered by inclusion with their improper
faces, form a regular lattice polytope of one higher rank, with vertex
figures K and automorphism group

\[
\operatorname{Aut}(2^K)=V\rtimes\operatorname{Aut}(K).
\tag{M1}
\]

The coordinate action is the natural vertex permutation action. This
construction theorem is a substantial retained input; neither its
polytopality nor the identification of the full group in (M1) is
inferred solely from an order calculation.

\textbf{Corollary C.Parity.} Halving \(2^K\) at its vertex endpoint
gives the exact group

\[
E_m\rtimes\operatorname{Aut}(K),\qquad
E_m=\{x\in\mathbf F_2^m:\sum x_i=0\},
\tag{M2}
\]

of order \(2^{m-1}|\operatorname{Aut}(K)|\), acting as the full
type-preserving group of the regular hypertope in Corollary C.Halving.

\textbf{Proof.} The vertex graph of \(2^K\) is the m-cube, bipartitioned
by coordinate parity. Its color-preserving subgroup of (M1) is exactly
(M2), so Corollary C.Halving identifies the whole halving group.
Alternatively, let a flip the coordinate of a base vertex v of K and let
b be its first distinguished reflection, carrying v to an adjacent
vertex \(v'\). The product \(a'b=(aba)b\) is translation by
\(e_v+e_{v'}\). The unchanged reflections generate Aut(K). Their
conjugates give every edge difference in the connected vertex graph of
K; along paths, these span \(E_m\). This independently proves generation
of the displayed subgroup. ∎

N013 uses the dual convention \(\widehat P=(2^{K^*})^*\). Consequently m
in its formulas is the number of facets of K, and reversing all types
converts (M2) exactly to its construction.

\textbf{Corollary C.Polygon (polygon inputs).} For a p-gon K,
\(p\ge 3\), the result is a self-dual orientable regular map of type
\{p,p\}, with

\[
V=F=2^{p-1},\qquad E=p2^{p-2},\qquad
|\operatorname{Aut}|=p2^p,\qquad
g=1+2^{p-3}(p-4).
\tag{M3}
\]

\textbf{Proof.} The two vertex colors give V and F after ordering the
map types as (color 0, old square, color 1). Each edge of the p-gon
specifies \(2^{p-2}\) coordinate squares, giving E. The fork diagram
becomes \{p,p\}. A single coordinate flip conjugates the two new
endpoint reflections and exchanges the colors, supplying the duality.
Extend the reflection-sign character of the dihedral Aut(K) trivially
over \(E_m\). Each of the three new distinguished reflections has sign
\(-1\), so the chamber orientation is well defined and reverses across
every adjacency. This proves orientability. Euler's formula
\(V-E+F=2-2g\) gives (M3). ∎

For \(p=3\) the map is the tetrahedron: the group is
\(E_{3}\rtimes S_{3}\) of order 24 with the spherical \(A_{3}\) marking
and its universal Coxeter geometry. For \(p=4\) it is the square toroid
conventionally denoted \(\{4,4\}_{2,2}\). Section 6 proves that lattice
identification and identifies the original dual generalized cube with
\(\{4,4\}_{4,0}\) by explicit face-poset and marking maps.

\textbf{Corollary C.Spherical (all locally spherical inputs and
exceptional orders).} In output rank \(n\ge 4\), the locally spherical
cases of N013's dual generalized-cube construction are obtained from K
equal to the \((n-1)\)-simplex or \((n-1)-cube\), and additionally the
dodecahedron for \(n=4\) and the 120-cell for \(n=5\). Their types are
respectively spherical \(D_n\), affine \(\widetilde B_{n-1}\), and the
rank-four/rank-five compact-hyperbolic forks. The last two finite groups
have orders

\[
2^{11}\cdot120=245760,
\qquad 2^{119}\cdot14400.
\tag{M4}
\]

\textbf{Proof.} Put \(L=K^*\) to use the coordinate model. Each of its
two new vertex residues is L itself. If the output is locally spherical,
L must therefore be a universal spherical polytope. Retain the classical
finite string-Coxeter list. The lattice hypothesis excludes digonal
rank-two sections, so every relevant adjacent label is at least three.
In rank at least four a proper three-node fork with equal labels q,q is
spherical only for \(q=3\), where q is the first label of L. In rank
five, the remaining candidate \(F_{4}=\{3,4,3\}\) is excluded because
its proper four-node fork is affine. The finite list leaves exactly the
simplex, cross-polytope, icosahedron and 600-cell, the last two only in
their respective ranks. For rank at least six only the first two remain.

Conversely, all proper faces of these L are simplices. At an old j-face
of \(2^L\), \(j\ge 2\), let F be its associated \((j-1)\)-face of L. The
lower and upper sections vary independently. After halving, the maximal
residue is the product of the compressed cube \(2^F\) with the spherical
section L/F. As F has j vertices, the first factor has group
\(E_j\rtimes S_j\) and its \(D_j\) marking; its order \(2^{j-1}j!\)
equals that of the finite Coxeter group \(D_j\), proving equality with
the universal spherical geometry. Here \(D_{2}=A_{1}\times A_{1}\) and
\(D_{3}=A_{3}\). Vertex residues are L, so all maximal residues, and
consequently all proper residues, are genuinely spherical geometries. A
diagram alone is not used as a substitute for those residues.

The new global diagrams follow from Corollary C.Halving. The retained
finite, affine and compact-hyperbolic Coxeter classifications identify
them as stated. Formula (M2), with \(m=12\) or 120 and the classical
orders \(|\operatorname{Aut}(dodecahedron)|=120\) and
\(|\operatorname{Aut}(120-cell)|=14400\), gives (M4). For the simplex
\(m=n\) and \(|\operatorname{Aut}(K)|=n!\), giving \(D_n\) order
\(2^{n-1}n!\). For the \((n-1)-cube\), \(m=2(n-1)\) and
\(|\operatorname{Aut}(K)|=2^{n-1}(n-1)!\), giving \(2^{3n-4}(n-1)!\).
The finite hyperbolic examples are proper hypertopes because the fork
diagrams have a vertex of degree three. ∎

The exact all-rank toroidal labels require more than these orders and
residues. They are established in the next section by an explicit
coordinate identification.

\subsection{6. Exact toroidal face posets and translation
lattices}\label{exact-toroidal-face-posets-and-translation-lattices}

Let \(r\ge 2\) and let L be the r-dimensional cross-polytope with
vertices labelled \(+e_i,-e_i\), \(1\le i\le r\). A proper face of L
chooses at most one vertex from each antipodal pair. Write \(Q_r\) for
the cubical tessellation of \(\mathbb R^r\) whose vertices are
\(\mathbb Z^r\) and whose cells are products of unit intervals and
integer singletons. The quotient \(Q_r/4\mathbb Z^r\) has a formal least
and greatest face adjoined when viewed as an abstract polytope. These
formal faces are also understood below.

For each pair of binary coordinates belonging to \(+e_i,-e_i\), use the
Gray map

\[
(0,0)\mapsto0,\quad(1,0)\mapsto1,\quad(1,1)\mapsto2,\quad(0,1)\mapsto3
\quad\text{in }\mathbb Z/4\mathbb Z.
\tag{T1}
\]

\textbf{Lemma C.Gray.} The product of these maps is a rank-preserving
face-poset isomorphism \(2^L\) ≅ \(Q_r/4\mathbb Z^r\). Binary weight
modulo two becomes \(\sum _i\) \(x_i\) modulo two in the cubical
coordinates.

\textbf{Proof.} In Danzer's construction a proper face of \(2^L\) allows
the bits indexed by a face of L to vary, and fixes every other bit. Thus
in any antipodal pair either both bits are fixed or exactly one bit
varies. Under (T1), varying a single bit describes precisely one edge of
the four-cycle: its four possible edges arise from the two choices of
varying bit and the two values of the fixed bit. A face therefore
becomes a product of four-cycle vertices and edges. Conversely every
such product uniquely specifies the varying bit in each free pair and
the remaining fixed bits. Its dimension is the number of free pairs.
Inclusion is preserved and reflected coordinatewise, and formal faces
match. These products are exactly the cubical cells modulo
\(4\mathbb Z^r\), with their incidence. The four cases of (T1) have
binary parities 0,1,0,1 respectively, proving the last assertion. ∎

This is a face-poset identification, rather than an inference from
orders. It also handles the dual in the target source explicitly. For a
cubical face F with free coordinate set S, write its free intervals as
\([m_i,m_i+1]\) for \(i\in S\) and its fixed coordinates as \(m_j\) for
\(j\notin S\). The dual face F\^{}* has fixed coordinate \(m_i+1/2\) for
\(i\in S\) and free interval \([m_j-1/2,m_j+1/2]\) for \(j\notin S\).
This reverses containment: whenever an endpoint is fixed in a subface of
F, the corresponding free interval of the dual contains the fixed
midpoint of F\^{}*. It is a bijection onto the shifted cubical
tessellation, with dimensions adding to r. Subtracting
\((1/2,\ldots ,1/2)\) identifies that tessellation with \(Q_r\). The
construction commutes with translations by \(4\mathbb Z^r\), and
consequently induces an order-reversing bijection on the quotient. Thus
for the r-cube \(K=L^*\), the source's polytope \((2^{K^*})^*\) is
indeed the cubical toroid \(\{4,3^{r-2},4\}_{(4,0,\ldots,0)}\). The
string is \{4,4\} when \(r=2\). Reversing the source marking moves its
last-end halving to first-end halving of \(2^L\). Type reversal is kept
explicit; no unproved compatibility of a chosen self-duality with
halving is needed.

\subsubsection{6.1. Markings and the affine
quotient}\label{markings-and-the-affine-quotient}

Use the following reflections of \(\mathbb R^r\):

\[
\begin{aligned}
a(x_1,x_2,\ldots,x_r)&=(1-x_1,x_2,\ldots,x_r),\\
b_i&:\ x_i\leftrightarrow x_{i+1}\quad(1\le i<r),\\
c(x_1,\ldots,x_{r-1},x_r)&=(x_1,\ldots,x_{r-1},-x_r).
\end{aligned}
\tag{T2}
\]

They generate \(W=\mathbb Z^r\rtimes B_r\), where \(B_r\) is the signed
permutation group. Indeed the \(b_i,c\) generate \(B_r\), including the
reflection \(x_1\mapsto -x_1\); its product with a is translation by
\(±e_1\). Signed permutations then supply every integer translation. The
simplex \(1/2\ge x_1\ge x_2\ge \cdots \ge x_r\ge 0\) is a fundamental
chamber. One way to verify this without a presentation theorem is to
translate each coordinate to \([-1/2,1/2]\), make it nonnegative, and
order the coordinates; a point with strict inequalities has a unique
such representative. Its reflected walls are the displayed generators.
Its vertices are the barycentres of a nested flag of cubical faces, so
its chamber system is the barycentric subdivision of \(Q_r\).

The marking in (T2) also agrees with Danzer's marking under (T1).
Flipping the \(+e_1\) bit induces \(x_1\mapsto 1-x_1\) modulo four;
interchanging antipodal pairs induces the \(b_i\); swapping the two bits
in the last pair induces \(x_r\mapsto -x_r\). These are the usual base
vertex/base flag generators of \(2^L\). Hence the isomorphism of Lemma
C.Gray identifies the marked action, not only the unmarked face poset.
The action of W on \(Q_r/4\mathbb Z^r\) has kernel exactly
\(\Lambda =4\mathbb Z^r\): an affine signed permutation acting trivially
on all vertices has translation zero modulo four, and its linear part is
the identity modulo four; testing the unit vectors forces that signed
permutation to be the identity.

Put \(b=b_1\) and replace a by

\[
d=aba:\ (x_1,x_2,x_3,\ldots,x_r)\mapsto(1-x_2,1-x_1,x_3,\ldots,x_r).
\tag{T3}
\]

Let \(E=\{v\in \mathbb Z^r\) : \(\sum v_i\) is even\}.

\textbf{Lemma C.Affine.} \(H=\langle d,b_1,\ldots ,b_{r-1},c\rangle\) is
exactly \(E\rtimes B_r\). It has index two in W and contains
\(\Lambda\). Its fundamental chamber is

\[
x_1\ge x_2\ge\cdots\ge x_r\ge0,\qquad x_1+x_2\le1.
\tag{T4}
\]

Its chamber system is the halved cubical geometry, and the marked finite
halving group and geometry are \(H/\Lambda\) and the quotient of that
chamber system by \(\Lambda\), respectively.

\textbf{Proof.} Every displayed generator preserves the parity of an
integral vertex's coordinate sum, so \(H\le E\rtimes B_r\). It contains
\(B_r\). The product of d with the \(B_r-reflection\)
\((x_1,x_2)\mapsto (-x_2,-x_1)\) is translation by \(e_1+e_2\).
Conjugating by \(B_r\) gives every \(±e_i±e_j\). These span E: subtract
multiples of \(e_i-e_r\), i\textless r, to leave an even multiple of
\(e_r\), and use \(2e_r=(e_r+e_1)+(e_r-e_1)\). This works also for
\(r=2\). Thus equality holds and the index is two. Since every \(4e_i\)
is in E, \(\Lambda \le H\).

The region (T4) is the union of the preceding W chamber and its
reflection in the a-wall \(x_1=1/2\). On their common wall the two
chambers glue; the other wall \(x_1=x_2\) is paired with its a-image
\(x_1+x_2=1\). All other walls agree. Because \(W=H\sqcup Ha\) and the
interiors of W chambers are disjoint, these doubled chambers tile
\(\mathbb R^r\) under H, each once. Their wall reflections are exactly
\(d,b_1,\ldots ,b_{r-1},c\). Equivalently, their chamber labels are the
old cubical flags with the initial vertex chosen in one parity class,
pairing the two old flags differing only at that vertex. The two ends of
the old edge become the two new vertex types; the higher faces retain
their type. This is precisely the graph-halving construction of Section
4.

Finally, translations in \(\Lambda\) preserve both vertex colors and
take cubical faces to cubical faces. The same flag pairing therefore
commutes with \(\Lambda\). Every finite chamber has a lift, and two
lifts yield the same finite chamber exactly when they differ by
\(\Lambda\). Wall adjacencies are the marked reflections modulo
\(\Lambda\). The maximal parabolics are the images of the corresponding
wall subgroups, so the typed coset geometry also agrees. Regularity,
thinness and residual connectedness of the finite geometry follow from
the already proved halving theorem; they are not inferred solely from
being a group quotient. ∎

For \(r\ge 3\), the walls (T4) have the affine fork diagram
\(\widetilde B_r\): d and \(b_1\) commute; both meet \(b_2\) with
product order three, successive \(b_i\) have order three, and
\(b_{r-1}c\) has order four. All other distinct wall pairs commute. This
follows directly from (T2)--(T3), or from their mirror angles. Thus its
rank-(r+1) quotient has the source's type \(\{4,3^{r-3},3/3\}\). Here
the quotient lattice is exactly \(4\mathbb Z^r\) in the original
coordinates; the subscript \((4,0,\ldots ,0)\) denotes the lattice
generated by that vector's signed-permutation orbit. This proves the
source's all-rank toroidal identification, including the marking. Its
group order is \([E:4\mathbb Z^r]|B_r|=(4^r/2)2^r\) \(r!=2^{3r-1}r!\),
agreeing with the separate parity-module calculation. The order is a
check, not the proof.

The coordinate convention can be checked directly against Ens's primary
paper: its equations (3.1), (4.1) are (T2), (T3) for \(r=3\); its §2
definition of \(\Lambda _s\) uses the signed-permutation orbit of s, and
Theorem 4.3 denotes the corresponding Coxeter-complex quotient by that
lattice. Thus \(s=(4,0,0)\) here is literally the same parameter. The
target extends this convention to larger r. No classification theorem
about all invariant lattices is needed. See {[}Ens, pp.70--75{]}. Only
that lattice convention is borrowed here, not its classification
theorem.

\subsubsection{6.2. The rank-three square
map}\label{the-rank-three-square-map}

When \(r=2\), reorder the halved types as (even vertex, old square, odd
vertex). An old unit square joins its two even corners as an edge of the
new map; the four even neighbors of an old odd vertex bound a new square
face. This describes the ordinary square tessellation with vertex
lattice E and edge directions

\[
f_1=(1,1),\qquad f_2=(-1,1).
\tag{T5}
\]

Every even vertex has four such incident edges, every displayed
quadrilateral is a face, and the incidences are exactly those of
halving. Thus this is an isomorphism of maps, not a comparison of their
counts. In this basis

\[
4e_1=2f_1-2f_2,\qquad4e_2=2f_1+2f_2.
\]

The quotient lattice has basis \((2,-2),(2,2)\) in square coordinates.
Negating the first vector gives the customary pair \((-2,2),(2,2)\), so
the map is \(\{4,4\}_{2,2}\). Lemma C.Gray and the explicit duality
already identify the original map as \(\{4,4\}_{4,0}\). This proves both
\(p=4\) identifications in {[}MW{]}.

\subsection{7. Six Janko-group geometries and their complete
census}\label{six-janko-group-geometries-and-their-complete-census}

The second application takes the odd branch of Corollary C.Odd. Retain
the GAP TomLib J1 permutation model, of order 175560. The construction
certificate constructs a marking (a,b,c,d) within that model with

\[
|ab|,|ac|,|ad|,|bc|,|bd|,|cd|=(5,2,2,3,2,5),
\] \[
(abc)^5=((dcb)^5a)^3=1.
\tag{J1}
\]

It also completes a Todd--Coxeter calculation for the literal
presentation with the four involution relations and (J1): the index over
\(\langle b,c,d\rangle\) is 1463. That subgroup is a quotient of the
finite Coxeter group {[}3,5{]} of order 120, so the presented group has
order at most \(1463·120=175560\). The found generating tuple in the
TomLib model supplies a quotient of that order; hence this is the source
presentation of J1. A valid permutation action alone would not prove the
upper bound: the completed exact coset enumeration is a charged
computational input.

Set \(w=(bcd)^5\). The certificate proves that w is a nonidentity
central involution of \(\langle b,c,d\rangle\). The six source-defined
markings and their exact invariants are:

{\def\LTcaptype{none} % do not increment counter
\begin{longtable}[]{@{}
  >{\raggedright\arraybackslash}p{(\linewidth - 6\tabcolsep) * \real{0.2500}}
  >{\raggedright\arraybackslash}p{(\linewidth - 6\tabcolsep) * \real{0.2500}}
  >{\raggedright\arraybackslash}p{(\linewidth - 6\tabcolsep) * \real{0.2500}}
  >{\raggedright\arraybackslash}p{(\linewidth - 6\tabcolsep) * \real{0.2500}}@{}}
\toprule\noalign{}
\begin{minipage}[b]{\linewidth}\raggedright
Geometry
\end{minipage} & \begin{minipage}[b]{\linewidth}\raggedright
Marking
\end{minipage} & \begin{minipage}[b]{\linewidth}\raggedright
Pair orders ab,ac,ad,bc,bd,cd in that marking
\end{minipage} & \begin{minipage}[b]{\linewidth}\raggedright
Maximal-parabolic orders
\end{minipage} \\
\midrule\noalign{}
\endhead
\bottomrule\noalign{}
\endlastfoot
\(\Gamma 1\) & (a,b,c,d) & 5,2,2,3,2,5 & 120,20,20,60 \\
\(\Gamma 2\) & (a,wb,c,d) & 5,2,2,6,2,5 & 120,20,20,660 \\
\(\Gamma 3\) & \((a,b,c,c^d)\) & 5,2,2,3,3,5 & 120,20,60,60 \\
\(\Gamma 4\) & \((b^a,b,c,d)\) & 5,3,2,3,2,5 & 120,120,20,60 \\
\(\Gamma 5\) & \((b^a,b,c,c^d)\) & 5,3,3,3,3,5 & 120,120,60,60 \\
\(\Gamma 6\) & \(((wb)^a,wb,c,d)\) & 5,6,2,6,2,5 & 120,120,20,660 \\
\end{longtable}
}

Order-120 parabolics are \(C_{2}\times A_{5}\), order-20 ones are
\(D_{20}=C_{2}\times D_{10}\), order-60 ones A5 and order-660 ones
\(\operatorname{PSL}_{2}(11)\); D10 here has order ten. Structure
identifications use the retained exact GAP algorithms and named group
models.

\textbf{Corollary C.J1.} Each of these six marked coset systems is a
thin regular residually connected rank-four geometry with full
type-preserving group J1. They are pairwise nonisomorphic, even allowing
type permutations. These are L088 Theorems 3.2, 4.2, 5.2, 5.4, 6.2 and
7.2, including their generation lemmas.

\textbf{Proof.} Lemma C.Twist gives the central twist
\(\Gamma 1\to \Gamma 2\) and preserves its first parabolic. Every
endpoint operation has odd endpoint order five, so Corollary C.Odd
recovers the same group. The construction certificate verifies the two
actual subgroup equalities

\[
\langle b,c,d\rangle\cap\langle b,c,d\rangle^a=\langle c,d\rangle,
\qquad
\langle a,b,c\rangle\cap\langle a,b,c\rangle^d=\langle a,b\rangle.
\tag{J2}
\]

They have order ten. The second gives \(\Gamma 3\), the first
\(\Gamma 4\). In \(\Gamma 3\),
\(\langle b,c,c^d\rangle =\langle b,c,d\rangle\) because \(|cd|=5\), and
\(\langle c,c^d\rangle =\langle c,d\rangle\); hence the same first
equality gives \(\Gamma 5\). In \(\Gamma 2\) its first parabolic is
unchanged by the twist, so that equality gives \(\Gamma 6\). Thus after
the seed and twist cases, four IP proofs reduce to the two meets (J2).
The construction certificate additionally verifies all 256 ordered
marked-subset meets for each of the six markings, so the finite result
does not depend on a mistaken reuse of a subgroup identity.

For each marking it checks all sixteen equalities (C3), without
presuming that a C-group with a nonstring diagram is a geometry. Lemma
C.HellyCount and U.Geometry now give all the geometric conclusions. The
six multisets of maximal-parabolic orders in the table are distinct, so
no type-permuting isomorphism identifies two geometries. ∎

The definition of \(\Gamma 3\) in L088 §5 duplicates an A5 parabolic,
whereas its printed figure labelled \(\Gamma 3\) duplicates
\(C_{2}\times A_{5}\); the definition of \(\Gamma 4\) has the reverse
discrepancy. The source figure labels \(\Gamma 3\) and \(\Gamma 4\) are
therefore transposed relative to its definitions. The table follows the
definitions. This is a source correction, not a seventh construction.

The two source-mentioned failures also follow from the same certificate.
The markings \((a,wb,c,c^d)\) and (a,b,c,wd) still generate J1 and
satisfy typed Helly, but each has eight failing ordered IP meets; their
panel orders are respectively (2,4,2,2) and (2,2,2,4). Therefore they
are not thin. These failures demonstrate why neither endpoint generation
nor flag transitivity is a substitute for IP.

\subsubsection{7.1. Exhaustive proper-residue
extension}\label{exhaustive-proper-residue-extension}

The original geometry census can be replaced by a short complete search
on proper residues. The required group data are the complete TomLib
table of J1 subgroup-conjugacy classes, not a list of its geometries.

\textbf{Theorem C.Census.} Let G be finite and let \(M_{1},\ldots ,M_s\)
represent all conjugacy classes of maximal proper subgroups. For each
\(M_j\) enumerate, up to \(M_j-conjugation\), all unordered k-element
sets of involutions whose markings satisfy IP and typed Helly. Append
every involution of G outside the subgroup generated by each retained
set and apply the same tests. This finds every rank-(k+1) thin regular
residually connected geometry whose full type-preserving group is a
subgroup of G. To restrict to group G, retain exactly the extensions
generating G.

\textbf{Proof.} For any such rank-(k+1) geometry, its distinguished
marking consists of independent involutions. Omitting one gives a proper
subgroup P of its full group and hence of G, so P lies in a conjugate of
one \(M_j\). Its marked coset geometry is the corresponding rank-k
residue, thus satisfies both IP and typed Helly. Conjugation into
\(M_j\) and then within \(M_j\) puts its unordered markers in the
enumerated list. Its omitted involution lies outside P and occurs in the
extension loop. Conversely, the tested IP and Helly conditions give
exactly the asserted geometries by U.Geometry. More formally, for each
\(M_j\) take E to be its admissible unordered k-marker sets and Z to be
pairs (T,t) with \(T\in E\), t an involution outside
\(\langle T\rangle\), and the required extension properties. Conjugation
by \(M_j\) acts on both sets, and \((T,t)\mapsto T\) is equivariant.
U.ExtensionOrbits identifies the extension orbits over a representative
T with the stabilizer orbits in its fiber. The certificate tries the
whole fiber, so it certainly meets each such orbit; its final conjugacy
keys remove duplicate extensions. Ordering the marker set merely orders
the type labels. ∎

For completeness, the residue statement used here requires no extension
of arbitrary residue automorphisms to G. The base-flag stabilizer acts
transitively and freely on the residue's chambers, and the residue is
thin and connected. Its full type-preserving group also acts freely, so
the stabilizer already is that full group. Walking inside the residue
identifies its panel generators. Thus local Helly is a necessary
condition for these genuine residues, not an arbitrary extra filter.

\textbf{Corollary C.Classification.} Up to isomorphism allowing a
permutation of types, the six geometries in Corollary C.J1 are all the
thin regular residually connected rank-four geometries with group J1.
There are none of rank at least five.

\textbf{Proof.} The exhaustive certificate executes Theorem C.Census for
\(k=3\) and \(k=4\). The retained complete table has forty subgroup
classes. Its nonzero containment entries determine seven maximal classes
of orders 660,168,120,114,110,60,42. The local counts are:

{\def\LTcaptype{none} % do not increment counter
\begin{longtable}[]{@{}
  >{\raggedleft\arraybackslash}p{(\linewidth - 10\tabcolsep) * \real{0.1667}}
  >{\raggedleft\arraybackslash}p{(\linewidth - 10\tabcolsep) * \real{0.1667}}
  >{\raggedleft\arraybackslash}p{(\linewidth - 10\tabcolsep) * \real{0.1667}}
  >{\raggedleft\arraybackslash}p{(\linewidth - 10\tabcolsep) * \real{0.1667}}
  >{\raggedleft\arraybackslash}p{(\linewidth - 10\tabcolsep) * \real{0.1667}}
  >{\raggedleft\arraybackslash}p{(\linewidth - 10\tabcolsep) * \real{0.1667}}@{}}
\toprule\noalign{}
\begin{minipage}[b]{\linewidth}\raggedleft
Maximal subgroup order
\end{minipage} & \begin{minipage}[b]{\linewidth}\raggedleft
Involutions
\end{minipage} & \begin{minipage}[b]{\linewidth}\raggedleft
Triple orbits
\end{minipage} & \begin{minipage}[b]{\linewidth}\raggedleft
Quadruple orbits
\end{minipage} & \begin{minipage}[b]{\linewidth}\raggedleft
Regular triple orbits
\end{minipage} & \begin{minipage}[b]{\linewidth}\raggedleft
Regular quadruple orbits
\end{minipage} \\
\midrule\noalign{}
\endhead
\bottomrule\noalign{}
\endlastfoot
660 & 55 & 60 & 588 & 42 & 6 \\
168 & 7 & 3 & 3 & 2 & 0 \\
120 & 31 & 108 & 616 & 45 & 7 \\
114 & 19 & 12 & 42 & 0 & 0 \\
110 & 11 & 2 & 4 & 0 & 0 \\
60 & 23 & 58 & 221 & 6 & 2 \\
42 & 7 & 2 & 2 & 0 & 0 \\
Total & --- & 245 & 1476 & 95 & 15 \\
\end{longtable}
}

The finite conjugation action is used here without changing its orbit
relation. Conjugation on the actual set of involutions has a permutation
image; every element of that image has an ambient conjugating lift, so
its induced subset orbits are exactly the needed conjugacy orbits.
Faithfulness of this action is not assumed. No local generated subgroup
is required to equal its containing maximal subgroup. The 95 retained
triples reduce to 70 J1-conjugacy classes. All 100016 possible
outside-involution extensions are tested; forty pass with generated
group J1, collapsing to precisely six classes. Their maximal-parabolic
order multisets are the six distinct ones in Corollary C.J1, which
already supplies six pairwise nonisomorphic constructions.

Here is the exact conjugacy test, so equality of numerical invariants is
not mistaken for isomorphism. There is one class of 1463 involutions in
J1, with centralizer order 120. Fix \(a_{0}\). For each x in an
unordered marker set T choose \(g_x\) with \(x^{g_x}=a_{0}\), enumerate
\(T^{g_xc}\) for every \(c\in C_G(a_{0})\), and take the least sorted
tuple of involution indices. These are exactly the conjugates of T
containing \(a_{0}\): every conjugator carrying x to \(a_{0}\) differs
from \(g_x\) on the right by such a c.~Consequently keys agree exactly
for inner-conjugate unordered markings. There is no assumption that all
automorphisms of J1 are inner. The six different parabolic-order
multisets prevent any further isomorphism from merging those classes,
while the complete six-class upper bound prevents any additional
geometry.

For \(k=4\), the fifteen local regular quadruple classes remain fifteen
after J1-conjugacy deduplication. All 21352 extensions by an involution
outside their generated subgroup fail the full regular-geometry tests.
This stage does not discard proper-subgroup targets. If a regular
geometry of rank \(r\ge 5\) existed for J1, fixing \(r-5\) objects of a
base chamber would yield a rank-five regular residue with group
contained in J1; Theorem C.Census would have found it. Thus all higher
ranks are excluded. This does not assert nonexistence of C-group
markings that fail typed Helly. ∎

The source's successful central modifications in §8, when they generate
J1 and yield a thin regular geometry, are consequently covered by the
same exhaustive classification. The two displayed failures were checked
separately above. The construction certificate also verifies the
incidental source equality
\(N_G(\langle a,b,c\rangle )=\langle a,b,c\rangle\). Its initial vertex
figure is the universal spherical {[}3,5{]} geometry, since that
parabolic has the full Coxeter order 120; its facet is the antipodal
hemi-dodecahedral quotient of {[}5,3{]}, since the A5 quotient has order
60 and kernel the unique central order-two subgroup of
\(C_{2}\times A_{5}\). These conventional spherical identifications use
the named classical Coxeter models.

\subsection{8. Geometric Petrials and the exact faithful
image}\label{geometric-petrials-and-the-exact-faithful-image}

The operation in this section is that of Bonheure--Buekenhout--Leemans
(2001), with equal geometric Petrie sets identified. Its principal
theorem and valid numbered structural conclusions are the targets. The
unconditional embedding assertion is corrected below. No source
conclusion is used as a hypothesis, and neither a finite group nor
finite residues are assumed.

Let \(\Gamma\) be a thin regular residually connected rank-three
geometry. Label its types vertex, edge and face. At a base chamber
\(C=(v,e,f)\), let a change the face, b the edge and c the vertex. Lemma
C.Gallery gives

\[
 G=\operatorname{Aut}(\Gamma)=\langle a,b,c\rangle,
 \quad G_v=\langle a,b\rangle,\quad G_e=\langle a,c\rangle,
 \quad G_f=\langle b,c\rangle,\quad G_v\cap G_e=\langle a\rangle.
 \tag{P1}
\]

The other two pair intersections are respectively \(\langle b\rangle\)
and \(\langle c\rangle\), and the chamber stabilizer is trivial. Any two
of a,b,c generate a dihedral group, possibly infinite: a word in two
involutions alternates after cancellation and is a power of their
product or such a power times one involution. For distinct nonidentity
involutions x,y, the intersection
\(\langle xy\rangle \cap \langle x\rangle\) is trivial. Otherwise x is a
power of xy and commutes with it, forcing \(xy=yx\) and \(|xy|=2\); then
\(x=xy\) would force \(y=1\). This proves the rotation/reflection fact
used below without a finiteness assumption.

\subsubsection{8.1. Geometric operation and its faithful
image}\label{geometric-operation-and-its-faithful-image}

Put \(w=cab\). Starting at a chamber, change vertex, face, then edge; at
the base chamber the result is wC. At gC, successive chamber flips act
by right multiplication on g, so iteration gives \(gw^iC\). The Petrie
face at C is the actual ordered pair

\[
 P(C)=\left(\{w^iv:i\in\mathbb Z\},\{w^ie:i\in\mathbb Z\}\right).
 \tag{P2}
\]

Equal pairs are one face; the vertex--edge incidence induced on these
sets is understood. Define \(\Pi \Gamma\) by retaining old vertices,
edges and their incidence, and adjoining these distinct faces, with
incidence to a face given by membership. This is equivariant:
\(P(gC)=gP(C)\). In particular \(wP(C)=P(C)\), and

\[
 wv=cv,\qquad w^{-1}e=be.
 \tag{P3}
\]

Thus a Petrie face contains v and cv at e, and e and be at v. Each pair
is distinct because the corresponding old panel change is nontrivial.

\textbf{Lemma C.PetrialAction (connectivity and exact action).}
\(\Pi \Gamma\) is a residually connected geometry. Its natural
automorphism action \(\rho\) has kernel

\[
 \ker\rho=\operatorname{core}_G\langle a\rangle
 =\begin{cases}\langle a\rangle,&a\in Z(G),\\1,&a\notin Z(G).\end{cases}
 \tag{P4}
\]

\emph{Proof.} Every incident old vertex--edge pair extends through an
old chamber to its Petrie face. For an object belonging to a Petrie
face, a step of the defining walk supplies an incident object of the
other retained type. These facts extend every flag to a chamber.

Write \(v_i=w^iv\) and \(e_i=w^ie\). Equation (P3) says that \(e_i\)
joins \(v_i\) and \(v_{i+1}\). The walk therefore connects every object
in its face residue; additional induced incidences preserve
connectivity. In the old residue at v, replace every two-step path
\(e_{1}–f–e_{2}\) with distinct edges by \(e_{1}–P(v,e_{1},f)–e_{2}\).
Thinness and (P3) justify the replacement. It connects all retained
edges in the new vertex residue, and every new incident face meets one
of them. The analogous replacement \(v_{1}–f–v_{2}\) proves connectivity
of each edge residue. Finally an old path through a face can be replaced
by a path in its connected vertex--edge residue. Thus the old
vertex--edge truncation, and consequently the new geometry with every
face attached, is connected.

For the action apply U.Action with X the disjoint union of old vertices
and edges. They are retained in Y, and every new face is uniquely
encoded by its two actual subsets of X. U.Action's core formula and (P1)
give the core in (P4). A normal subgroup of \(\langle a\rangle\) is
either trivial or all of \(\langle a\rangle\). The latter is normal in G
exactly when its unique nonidentity element a is central. This proves
the exact formula and constructs the faithful image \(G/ker\rho\). ∎

This is the same retained-object theorem used for graph halving in §4.1,
but the smaller retained object set can have a nontrivial pointwise
kernel. It asserts neither chamber transitivity of \(\Pi \Gamma\) nor
injectivity of \(\rho\) without further proof.

For an ordinary map, the three successive changes move across an edge to
its other endpoint, change the incident face, and continue along that
face. These are exactly the alternating corner choices in the
conventional Petrie walk. This proves the walk-definition assertion of
source Lemma 3.1; it does not forbid repetitions or replace geometric
faces by labelled walks.

\subsubsection{8.2. The thin obstruction and all conditional
conclusions}\label{the-thin-obstruction-and-all-conditional-conclusions}

\textbf{Theorem C.Petrial (Petrial obstruction and recovery).} If
\(\Pi \Gamma\) is thin, then \(|ab|=2\) or \(|ac|=2\). The action
\(\rho\) is faithful and regular and is the full automorphism group of
\(\Pi \Gamma\). Applying \(\Pi\) again recovers \(\Gamma\) by an
isomorphism fixing all retained vertices and edges.

\emph{Proof.} Set \(P=P(C)\). Suppose first that \(bP\ne P\). Both P and
bP contain e and be at v, so the vertex residue of \(\Pi \Gamma\)
contains the four-cycle e--P--be--bP--e. A finite cycle in a connected
graph of degree two is the entire graph: its vertices already have their
two neighbours, so no path can leave it. Lemma C.PetrialAction therefore
says that precisely two old edges are incident with v. In the old thin
connected vertex residue both edges have degree two, and each incident
face meets both. It is also a four-cycle. Its regular dihedral group
\(\langle a,b\rangle\) has four chambers, so \(|ab|=2\).

Suppose instead that \(bP=P\). Since w fixes P, \(ca=wb\) fixes P and e.
By the general freeness statement of U.Geometry, the automorphism group
of the thin \(\Pi \Gamma\) acts freely on chambers. Its edge--face panel
stabilizer embeds in the permutation group on that panel's two
completions, so \(\rho (ca)^2=1\). Equation (P4) puts \((ca)^2\) in
\(\langle a\rangle\). The rotation/reflection intersection in
\(\langle a,c\rangle\) is trivial, hence \((ca)^2=1\). Since \(c\ne a\),
this gives \(|ac|=2\). This proves the obstruction without presuming
that b always fixes P or that the new action is regular.

Interchanging vertex and edge interchanges b,c and changes cab to
\(bac=w^{-1}\), preserving the same Petrie sets after that type
interchange. We may consequently assume \([a,c]=1\). Then
\(bwb=w^{-1}\), \(bv=v\) and \(be=w^{-1}e\), so \(bP=P\). The subgroup
\(\langle b,w\rangle =\langle b,ca\rangle\) is transitive on the
incident vertex--edge pairs of P: its elements give \((v_i,e_i)\) and
\((v_i,e_{i-1})\) for every i. These are distinct pairs, and thinness
makes them all pairs at \(v_i\). Every object of P occurs in the walk.
The original chamber-transitive G is transitive on the faces P(gC), so
its induced action on \(\Pi \Gamma\) is chamber-transitive. U.Action
identifies its image with the full regular automorphism group. The
preimage of the new panel stabilizing (v,e) is exactly
\(G_v\cap G_e=\langle a\rangle\). U.Action's panel-index formula forces
\(ker\rho =1\).

At (v,e,P), the new panel involutions are therefore a,b,ca: the last
fixes (e,P) and sends v to cv; b fixes (v,P) and changes e; a exchanges
the two faces at (v,e). The new generators generate the same group since
\(c=(ca)a\). Their geometric marking is established by this action, not
by assuming a homomorphism from the old marking. Their second Petrie
step is \((ca)ab=cb\). The cb-orbits of v and e exhaust exactly the old
face's vertex--edge shadow: use the alternating normal form in its
residue group \(\langle b,c\rangle\) from (P1).

Those old face shadows are distinct. Otherwise let \(f\ne f'\) have the
same shadow and choose (v,e) in it, moving the base chamber if needed.
Thinness gives \(f'=af\). Both \(f,f'\) meet e and be at v. Their
four-cycle exhausts the old vertex residue, as above, forcing
\([a,b]=1\). Together with \([a,c]=1\) this makes a central in G, and
(P4) contradicts \(ker\rho =1\). Thus old faces correspond bijectively
to the second Petrie faces, with precisely their original incidences.
This proves the required isomorphism. The vertex--edge-swapped case
gives the same conclusion. ∎

The commuting pair is the source's linear-diagram conclusion. Theorem
C.Petrial proves its Theorem 1 and the thin-only formulation in its
abstract, together with Lemma 3.2(iii),(iv). Lemma C.PetrialAction gives
Lemma 3.2(i). No finite-residue assumption occurs in the walk,
connected-cycle, dihedral or gallery arguments.

\subsubsection{8.3. Corrected assertions and exact
controls}\label{corrected-assertions-and-exact-controls}

The unconditional Lemma 3.2(ii) must be read as the natural homomorphism
with kernel (P4), not as an embedding. Indeed take the complete
tripartite geometry with two objects of each type. Its full
type-preserving group is \(C_{2}^{3}\), regularly permuting eight
chambers. Let a,b,c independently change the face, edge and vertex
coordinates. The element cab changes all three, so every Petrie face
consists of both vertices and both edges. Deduplication leaves one face.
Thus \(\Pi \Gamma\) is the complete tripartite geometry with type sizes
(2,2,1), its full group is \(C_{2}^{2}\), and the kernel is
\(\langle a\rangle\). An order-eight group cannot embed in this
order-four automorphism group. This also gives a commuting-pair diagram
with nonthin Petrial, so the principal implication has no converse in
this diagram convention.

The unnumbered geometric length/order equality also cannot be
unconditional. Take \(C_{2}\times S_{3}\), with a generating its central
\(C_{2}\) and b,c distinct transpositions generating \(S_{3}\). The
coset geometry with stabilizers in (P1) is a six-cycle on three vertices
and three edges, joined to two faces. It is thin, residually connected
and regular: its twelve chambers form the regular \(C_{2}\times S_{3}\)
action. On retained objects, cab acts as cb, a rotation of order three,
while cab itself has order six. Its geometric Petrie face has three
vertices and three edges, not length six. This assertion is unnecessary
for every proof above.

Both counterexamples are proved directly. A separately supplied optional
program checks seven finite examples using literal chamber changes and
deduplicated vertex--edge sets; it is not a premise for the theorem or
an exhaustive classification. The uniform proof above covers infinite
groups and residues. No source atlas or external map classification is
needed.

\subsection{9. Source conclusions and qualified
assertions}\label{source-conclusions-and-qualified-assertions}

The results of §§3--7 recover the six constructions and complete rank
classification of {[}HL{]}, and the halving theorem and principal
applications of {[}MW{]}. Section 8 recovers the principal thin-Petrial
theorem and valid structural assertions of {[}BBL{]}, together with the
corrected action statement. The exact editions are the ten-page 2005
journal article {[}HL{]}, the eleven-page first arXiv version of
{[}MW{]}, and the complete seven-page rendered authored text of
{[}BBL{]}. The last is not a claim of possession or collation of its
original PDF. No later publisher numbering is substituted for the
specified author version.

{\def\LTcaptype{none} % do not increment counter
\begin{longtable}[]{@{}
  >{\raggedright\arraybackslash}p{(\linewidth - 2\tabcolsep) * \real{0.5000}}
  >{\raggedright\arraybackslash}p{(\linewidth - 2\tabcolsep) * \real{0.5000}}@{}}
\toprule\noalign{}
\begin{minipage}[b]{\linewidth}\raggedright
Source conclusions
\end{minipage} & \begin{minipage}[b]{\linewidth}\raggedright
Derivation here
\end{minipage} \\
\midrule\noalign{}
\endhead
\bottomrule\noalign{}
\endlastfoot
{[}HL{]}, Theorems 3.2, 4.2, 5.2, 5.4, 6.2, 7.2 and their generation
lemmas & C.Endpoint, C.Odd, C.Twist and C.J1; the complete construction
certificate \\
{[}HL{]}, exhaustive rank-four classification and no higher ranks &
C.Census and C.Classification; the complete proper-residue extension
certificate \\
{[}MW{]}, Theorem 3.1, Proposition 3.2, Corollary 3.3 and index/diagram
assertions & C.Endpoint, C.Graph and C.Halving, with all lattice
hypotheses proved \\
{[}MW{]}, generalized-cube, polygon and locally spherical applications &
C.Parity, C.Polygon and C.Spherical \\
{[}MW{]}, original and halved toroidal names & C.Gray, C.Affine and the
explicit rank-three lattice basis change \\
{[}BBL{]}, Theorem 1, Lemma 3.1 and Lemma 3.2(i),(iii),(iv) &
C.PetrialAction, the walk-definition adapter and C.Petrial \\
{[}BBL{]}, unconditional Lemma 3.2(ii) and length/order assertion &
Exact action kernel and the two written counterexamples of §8.3 \\
\end{longtable}
}

There is one further rank qualification in {[}MW, Remark 4.1{]}. Retain
Pellicer's extension construction used there. Its endpoint-2s halving
has a three-node fork triangle with labels (p,p,s). In output rank at
least four and for finite \(p,s\ge 3\), this is a proper rank-three
diagram and is not finite, since \(2/p+1/s\leq1\). Hence the resulting
hypertope is not locally spherical. In rank three, however, all its
proper residues have rank at most two and the displayed finite
parameters give finite dihedral residues. The unqualified
non-local-sphericity wording does not apply in that rank under the
source's definition. This remark does not assert a new classification of
the extension construction.

The comparison in {[}MW{]} with every example in its reference {[}9{]},
and its unused alternating-semiregular identification, are not needed by
any proof here and are not asserted as rederived. The classical Coxeter
classifications and Danzer theorem are retained background; the actual
target toroidal identifications are proved in §6. The
\(\Gamma 3/\Gamma 4\) label correction in {[}HL{]} and the two Petrial
counterexamples are explicit mathematical qualifications, not additional
source papers or claims of historical correction priority.

\subsection{Appendix. Complete finite-certificate
interface}\label{appendix.-complete-finite-certificate-interface}

The written completeness proof is Theorem C.Census. Its finite inputs
and outputs are supplied by two complete GAP programs, distributed with
this manuscript:

{\def\LTcaptype{none} % do not increment counter
\begin{longtable}[]{@{}
  >{\raggedright\arraybackslash}p{(\linewidth - 4\tabcolsep) * \real{0.3000}}
  >{\raggedright\arraybackslash}p{(\linewidth - 4\tabcolsep) * \real{0.3000}}
  >{\raggedleft\arraybackslash}p{(\linewidth - 4\tabcolsep) * \real{0.4000}}@{}}
\toprule\noalign{}
\begin{minipage}[b]{\linewidth}\raggedright
Program
\end{minipage} & \begin{minipage}[b]{\linewidth}\raggedright
Exact mathematical role
\end{minipage} & \begin{minipage}[b]{\linewidth}\raggedleft
Bytes
\end{minipage} \\
\midrule\noalign{}
\endhead
\bottomrule\noalign{}
\endlastfoot
\texttt{phase40\_compress\_operations\_check.g} & Seed search; completed
literal-presentation index; all six markings; central twist and overlap
identities; all subgroup meets and typed Helly tests; both failed
operations & 4,458 \\
\texttt{phase40\_compress\_operations\_exhaustive.g} & Complete
maximal-subgroup list; local involution-set orbits; every admissible
extension; exact conjugacy keys; rank-four census and rank-five
exhaustion in every subgroup & 4,783 \\
\end{longtable}
}

These are the full executable proofs of the finite assertions, not
summaries or a lookup of the source's geometry list. They contain no
import of an older geometry census. The first program checks 256 ordered
marked-subset intersections and sixteen global Helly equalities per
successful marking. Its literal-presentation calculation completes a
Todd--Coxeter enumeration of index 1,463; the retained spherical
subgroup order 120 gives the full order bound. The second imports all
forty J1 subgroup classes from TomLib and derives the seven maximal
classes from their containment table. It tests all 100,016 rank-four
extensions and all 21,352 rank-five extensions specified in §7. Its six
rank-four class keys and empty rank-five output, together with the six
distinct maximal-parabolic order multisets, provide the exact claimed
isomorphism count.

The subgroup catalogue is a substantial retained finite-group input. Its
completeness, the J1 permutation model, and GAP's exact subgroup,
conjugacy, orbit and completed coset-enumeration algorithms are not
formalized or rederived here. The recorded independent execution used
GAP 4.15.1 and TomLib 1.2.11 and binds their actual data file. The
source geometry classification itself is not retained. The structural
checker binds the two unchanged programs and the recorded independent
output; its structural replay is not mislabelled as a fresh census
execution.

An explicit fresh execution is available through
\texttt{phase150\_module\_C\_run.py\ -\/-output-dir\ NEW}; it runs the
same two hash-bound programs in a temporary working directory and
records the installed data, all completion checks and exact class
certificate. The 6,337-byte runner is an additional charged executable
artifact. It refuses existing output paths and modified mathematical
programs.

The separate source ledger records all thirty-five principal, auxiliary
and corrected scope entries and their complete proof locations. The
shared U body, this entire manuscript, the two full programs and the
actual certificate runner/binding infrastructure are the written and
computational package. Separately priced optional controls do not
replace an argument. In particular, the Petrial finite examples
establish no infinite theorem. Source-length comparison requires a
separate matched measurement of the complete package; a page count for
this exposition alone is insufficient.

The existing Lean file \texttt{EndpointDoublingIntersections.lean}
proves only the fixed-overlap mixed-subgroup identity used in
C.Endpoint. Its theorem \texttt{mixed\_intersection\_identified} still
requires the application-specific conjugation, fixed subgroup and marked
intersections to be supplied. It does not formalize the full endpoint
criterion, graph halving, the J1 census, the lattice bridge or the
Petrial theorem. No complete source application in this module is
currently labelled Lean-verified. The independent finite certificates
and the written all-parameter proofs retain their separate roles.

\subsection{References}\label{references}

{[}U{]} \emph{Marked groups, kernels and intersections}, companion
manuscript, third edition. The precise interfaces used here are
U.Action, U.Geometry, U.Helly and U.ExtensionOrbits.

{[}BBL{]} D. Bonheure, F. Buekenhout and D. Leemans, \emph{On the
Petrials of thin rank 3 geometries}, Journal of Geometry \textbf{71}
(2001), 19--25, DOI 10.1007/s00022-001-8548-2. The entire authored
rendered text is the declared source edition; the corrected
unconditional assertions are discussed in §8.

{[}HL{]} M. I. Hartley and D. Leemans, \emph{On the thin regular
geometries of rank four for the Janko group J1}, Innovations in
Incidence Geometry \textbf{1} (2005), 181--190. The defining markings,
rather than the transposed figure labels, determine \(\Gamma 3\) and
\(\Gamma 4\).

{[}MW{]} A. Montero and A. Ivić Weiss, \emph{Locally Spherical
Hypertopes from Generalised Cubes}, arXiv:2001.11082v1, 29 January 2020.
This author version supplies the theorem numbering and application
contract used here.

{[}PT{]} C. A. Piedade and P. Tranchida, \emph{Constructing new
geometries: a generalized approach to halving for hypertopes},
arXiv:2405.19050v1. The geometric graph construction, the lattice
interpretation of nondegeneracy and the bipartiteness/index-two relation
are attributed to the halving literature. The precise parity
compatibility used in this manuscript is proved in §4; no broader
theorem is imported without its hypotheses.

{[}Danzer{]} L. Danzer, *Regular incidence-complexes and dimensionally
unbounded sequences of such. I*, North-Holland Mathematics Studies
\textbf{87} (1984), 115--127. Its generalized-cube construction and full
original semidirect-product automorphism group are retained in §5.

{[}MS{]} P. McMullen and E. Schulte, \emph{Abstract Regular Polytopes},
Cambridge University Press, 2002. The standard spherical Coxeter models,
convex-polytope group orders and usual polytope conventions are retained
background. The particular face-poset and toroidal identifications used
here are proved in §6.

{[}Ens{]} E. Ens, \emph{Rank 4 toroidal hypertopes}, Ars Mathematica
Contemporanea \textbf{15} (2018), 67--79; selected pp.70--75, especially
equations (3.1), (4.1) and the translation-lattice convention preceding
Theorem 4.3. These selected passages fix the notation; the
classification theorem is not an input to the specific all-rank quotient
proof here.

The remaining named inputs are the classical finite, affine and
compact-hyperbolic Coxeter classifications; standard orientation and
Euler-characteristic facts for maps; Pellicer's extension construction
only in the qualified remark of §9; and the exact GAP/TomLib group data
and algorithms specified in the certificate interface. No historical
novelty, general open-conjecture resolution or complete formal
verification is claimed by this assembly.

\end{document}
